% NichidaiMasterExam_R6_1st_1
\begin{tcolorbox} %%R6_1st_1
	$\fbox{1}$ 次の問いに答えよ。
\begin{enumerate}
	\item $V$を実数体$\R$上の線形空間とする。$V$の空でない部分集合$W$が$V$の部分空間であることの定義を書け。
	\item 実数体$\R$上の$n$次元ベクトル空間$\R^n$に対して、次の集合は部分空間となるかどうか理由を付けて判定せよ。ただし$\bm{v}\in \R^n$の第$i$成分を$v_i$として表す。
\begin{enumerate_roman}
	\item $W_1=\{\bm{v}\in \R^n; v_1+ v_2+ \cdots+ v_n= 0\}$
	\item $W_2= \{\bm{v}\in \R^n; v_1\times v_2\times \cdots\times v_n= 0\}$
\end{enumerate_roman}

	\item 定義域を$\R$とする実数値関数$f\colon \R\to \R$の集合$V$とする。$V$は実数体$\R$の線形空間である(これは示さなくてよい)。次の集合は$V$の部分空間になるかどうか理由をつけて判定せよ。
\begin{enumerate_roman}
	\item $U_1=\{f\in V; fは有界\}$
	\item $U_2= \{f\in V; fは連続でない\}$
\end{enumerate_roman}
\end{enumerate}
\end{tcolorbox}

\begin{souanswer} %%R6_1st_1_ans
	$\fbox{1}$
 \begin{enumerate}
   \item 線形空間$V$の空でない部分集合$Wに対し、$W$が$V$の部分空間であるとは、次の条件を満たすことである。
  \begin{enumerate_alpha}
   \item 加法について閉じている：任意の$x, y\in W$に対して、$x+ y\in W$
   \item スカラー倍について閉じている：任意の$x\in W$および任意のスカラー$c\in \R$に対して、$cx\in W$
\end{enumerate_alpha}
   \item
\begin{enumerate_roman}
   \item $W_1= \{\v\in \R^n; v_1+ v_2+ \cdots+ v_n= 0}$
   部分集合である。\\
\begin{enumerate_alpha}
   \item 非零性：零ベクトル$\bm{0}= (0, 0, \ldots, 0)$について$0+ 0+ \cdots+ 0= 0$なので$\bm{0}\in W_1$であり、$W_1\ne \emptyset$
   \item 線型結合：任意の$\bm{u}= (u_1, u_2, \ldots, u_n),\ \bm{v}=(v_1, v_2, \ldots, v_n)\in W_1$および任意の$a, b\in \R$をとる。\\
   $a\bm{u}+ b\bm{v}$の各成分は$au_i+ bv_i$であるから、各成分の総和を計算すると、
\begin{equation*}
   \sum_{i= 1}^n (au_i+ bv_i)= a\sum_{i= 1}^n u_i+ b\sum_{i= 1}^n v_i
\end{equation*}
   ここで$\bm{u}, \bm{v}\in W_1$より$\sum_{i= 1}^n u_i= 0$かつ$\sum_{i= 1}^n v_i= 0$であるため、したがって$a\bm{u}+ b\bm{v}\in W_1$となり、$W_1$は加法およびスカラー倍について閉じている。
\end{enumerate_alpha}
   \item $W_2= \{v\in \R^n; v_1\times v_2\times \cdots\times v_n= 0\}$
   $n= 1$のとき線形空間であるが、$n\ge 2$のとき線形空間でない。\\
   加法について閉じていない反例が存在する。$\R^n$の基底ベクトルを$\bm{e}_1= (1, 0, \ldots, 0),\ \bm{e}_2= (0, 1, 0, \ldots, 0)$とする。
\begin{itemize}
   \item $\bm{e}_1$の成分の積は$1\times 0\times \cdots \times 0= 0$であるから、$\bm{e}_1\in W_2$
   \item $\bm{e}_2$の成分の積は$0\times 1\times 0\times \cdots \times 0= 0$であるから、$\bm{e}_2\in W_2$
   \item $\bm{e}_1$と$\bm{e}_2$の成分の和は$(1, 1, 0, \ldots, 0)$であり、$n= 2$の場合の積は$1\times 1= 1\ne 0$となり、$W_2$に属さない。同様に$n> 2$の場合も$W_2$は加法に関して閉じないために部分空間ではない。
\end{itemize}
\end{enumerate_roman}
   \item
\begin{enumerate_roman}
   \item $U_1= \{f\in V; fは有界\}$
   部分空間である。
\begin{enumerate_alpha}
   \item 非空性：恒等的に0である零関数O(x)= 0$はすべての$x\in \R$に対して、$|O(x)|\le 0$を満たすため有界である。よって$0\in U_1$
   \item 線型結合についての閉性：任意の$f, g\in U_1$および任意の$a, b\in \R$をとる。
\begin{equation*}
   |f(x)|\le M_f,\ \ |g(x)|\le M_g
\end{equation*}
   が成り立つ。三角不等式より任意の$x\in \R$に対して、
\begin{align*}
   |(af+ bg)(x)|
   &= |af(x)+ bg(x)|\\
   &\le |a||f(x)|+ |b||g(x)|\\
   &\le |a|M_f+ |b|M_g
\end{align*}
   ここで$M= |a|M_f+ |b|M_g$とおくと、$M$は$x$に依存しない定数であるため、$af+bg$も有界である。すなわち$af+ bg\in U_1$となり、$U_1$は加法およびスカラー倍について閉じている。
\end{enumerate_alpha}
   \item $U_2= \{f\in V; fは連続でない\}$
      部分空間でない。\\
      部分空間であるための必要条件「零元を含むこと」を満たさない。
\end{enumerate_roman}
\end{enumerate}
\end{souanswer}
